\scp{Reviewing CEMP}{15-03}
Like CMP,
we find no support for CEMP in the historical phonology or the
morphosyntax from the available data in our Wallacean subgroups,
nor any that brings unity to the Wallacean subgroups.

\ili{As} discussed in detail in \srf{13-01-05},
the basis for claiming the development of PCEMP mid vowels *e
and *o is not supported by the data in our target region.
We also showed that most of the original claims upon which those
PCEMP mid vowels are based have since been withdrawn,
and others have been conceded as due to diffusion.
We noted the reconstruction of the mid vowels to a proto-language
appears to be more appropriate to \tsc{POc} than to PCEMP.
Mid vowels are also reconstructed in contiguous subgroups west
of the Blust line,
so whatever claims are made cannot be seen as an exclusive innovation.

In \srf{13-02-01}, we discussed \citeauthor{Blust1993}’s
(\citeyear[246]{Blust1993}) claims regarding
“the reduction of consonant clusters in the reflexes of reduplicated
monosyllables” as a phonological innovation distinguishing PCEMP
from PMP, and noted that it is problematic in many regards.
The reduction of medial *CC in languages that have no medial
clusters is to be expected.
But in PMP words shaped *CVC\sub{1}C\sub{2}VC,
in some cases C\sub{1} is retained in Wallacean subgroups,
more often C\sub{2} is retained,
and in still others they coalesce (merge) to become something
else: C\sub{1}C\sub{2}  → C\sub{3}.
So to account for the data in the region,
while there is widespread reduction,
nevertheless both medial historical consonants must be reconstructed.
Many PCEMP reconstructions in which PMP *CVC\sub{1}C\sub{2}VC > PCEMP *CVC\sub{2}VC
do not reflect the full data in the region,
but only lower level nodes within subgroups, or individual languages.
Reduction of medial CC clusters is also widespread in languages
west of the Blust line such as with \ili{Malay},
so the pattern is not exclusive.

\citet[258]{Blust1993} saw “the use of proclitic subject markers
on the verb” as a “morphosyntactic innovation” in support of PCEMP.
In \srf{13-03-03}, we noted that forms that are cognate are retentions
of PMP forms, and other forms are not cognate with each other,
and often not even cognate with forms in different branches of
the same subgroup, thus not being useful for higher level subgrouping purposes.
And we noted that proclitic subject markers are also found west
of the Blust line,
with many of those forms also being retentions from PMP.

While we find no support for CMP or CEMP,
nevertheless the languages \ili{east} of the Blust line are clearly
different from languages west
and north of the Blust line in varying ways
and to varying degrees.
But it is not entirely clear what that collective difference might be.
On the basis of the evidence available to date it is not based
on the historical phonology or morphosyntax,
which are foundational requisites for subgrouping arguments.

Like the problems discussed in relation to the lexicon relating
to CMP, many of these same issues apply to PCEMP reconstructions.
We find many PCEMP reconstructions to be both problematic and unhelpful.
They often ignore many exceptions in the region that do not support such a reconstruction.
They may reflect possibilities for some select languages,
but ignore the fuller data that prohibit such a reconstruction in other languages.
For example, in \srf{10-02} it was noted that the PMP form *dahun
‘leaf’ with the medial consonant was required to account for
the \ili{Alune} vowel shift to \it{loi-ni} ‘leaf’,
whereas PCEMP **daun ignores \citeauthor{Collins1983}' (\citeyear[40--52]{Collins1983}) arguments
requiring the trace of PMP *h.
In most cases, the data in the region can be derived through
normal sound correspondences and normal processes for those languages
and those subgroups from the PMP form.\footnote{
		Another example is purported PCEMP **ruyuŋ `dugong'
		which does not account for the initial consonant
		of \ili{Helong} (TB, west Timor) \it{duiŋ} `dugong'.
		Instead, \ili{Helong} \it{duiŋ}
		is completely regular from PMP *duyuŋ.}

In these cases, the purported PCEMP reconstruction creates unhelpful
clutter that skews the focus away from what is really happening
in the languages in the region.
\ili{As} noted in \citet{GrimesCIP},
things like PMP *ma-qasin ‘salty’ > PCEMP **maqasin ‘salty,
brackish’ may inflate numbers,
but “do not identify innovations,
nor do they add to our knowledge about linguistic patterns or
the prehistory of the region.
They simply add clutter,
and collectively distract from seeing what might be insightful
and legitimate in the prehistory.” Reconstructions such as PCEMP
*qulu ‘head’ that are identical with PMP also do not add new information about innovations.
We already know that the sound inventory
and reconstructions accounting for the data in the region require
something almost identical to PMP \citep{Grimes1991a}.
So “reconstructions” of this sort merely give the impression
of substance to something that appears on closer inspection to be an illusion.

Methodologically for certain PCEMP reconstructions,
Blust determined the PCEMP form
and gloss for the data in the Wallacean region (and sometimes
even further west) based on the forms
and glosses of the Oceanic data far to the \ili{east}.
We have seen this problem with PCEMP **keRa(nŋ) ‘hawksbill turtle’
(corrected to PMP *kəRa(ŋ) ‘shell,
turtle shell’; see \srf{13-01-05}),
PCEMP **kandoRa ‘cuscus’ (corrected to **kVndoR(a) ‘cuscus’ or
possibly **kVnduR(a) for the \tsc{Eastern Islands} node of STB
and possibly SHWNG; see \srf{14-03-04}),
and PCEMP *mans(a,ə)r ‘bandicoot’ (corrected to **mantəR ‘cuscus’
with incrementally similar forms for several Wallacean subgroups
and \tsc{\ili{Muna}-Buton} that may or may not be related to \tsc{POc}
*mʷajar ‘bandicoot’; see \srf{14-04}).
Given Blust’s assumptions of the branching of PMP (see \frf{01-02}),
it is methodologically unsound to determine the form
and meaning of a node further back up the tree on the basis of
the form and meaning of a node further down the tree,
when there are clear discrepancies or discontinuities as we have
seen with these contested terms.